% ==============================================================================
% Fichier : front/01_avant_propos.tex : SAFI (SIS589)
% ==============================================================================

\chapter*{Avant-propos}
\addcontentsline{toc}{chapter}{Avant-propos}
\markboth{Avant-propos}{Avant-propos}

Ce livre est le support de l'unité d'enseignement \textbf{SIS589, Supervision, Analyse Forensique et Gestion des
Incidents}, du Master
Professionnel en Sécurité des Systèmes Informatiques et des Données de
l'Université de Garoua (Faculté des Sciences, Département de Mathématiques et
Informatiques).

\medskip
Il prolonge, sur le terrain opérationnel, les enseignements de sécurité et
d'audit~: là où l'on apprend ailleurs à apprécier les risques et à simuler les
attaques, SIS589 apprend à \textbf{détecter, investiguer et maîtriser} ce qui
survient en temps réel. Le fil conducteur est unique~: une entreprise
fictive, \textbf{LiZbiZ}, transporteur de colis présent à Garoua, subit un
rançongiciel. Le cours suit l'incident sur trois temps (la supervision qui
aurait dû le voir venir, l'analyse forensique qui l'établit, la réponse qui
le maîtrise), à travers trois opérations situées dans le Nord~: \textbf{BÉNOUÉ},
\textbf{FARO} et \textbf{MAYO}.

\medskip
Chaque notion est ancrée à la fois dans les standards internationaux (ISO/IEC
27035 et 27037, NIST SP 800-86, RFC 3227, MITRE ATT\&CK) et dans le droit
camerounais et africain (lois n°~2010/012, 2010/013 et 2024/017, Convention
de Malabo, Convention de Budapest). Un fil théorique parcourt le livre~: le
\textbf{Trilemme CRO} (Confidentialité, Fiabilité, Opposabilité), par lequel
se juge toute trace et tout acte d'enquête.

\medskip
Les cas sont fictifs~: les noms, les dates et les montants le sont~; les
mécanismes techniques et les références normatives sont exacts. Le livre
s'accompagne de l'environnement \textbf{NDAP SUKULU}, qui fait pratiquer
chaque chapitre.

\bigskip
\noindent \IfFileExists{../../Logos/signature_avant_propos.png}{\includegraphics[height=1.6cm]{../../Logos/signature_avant_propos.png}}{\textit{MINKA MI NGUIDJOÏ Thierry Emmanuel}}
